% part3_domain2.tex — ส่วนที่ 3 ผลการดำเนินงาน ด้านที่ 2 ตัวชี้วัดที่ 1
% โครงสร้างนี้ทำซ้ำได้กับทุกด้าน/ทุกตัวชี้วัดของแบบประเมิน PA 4
% (คัดลอกไฟล์นี้เป็นต้นแบบสำหรับด้าน/ตัวชี้วัดอื่น ๆ แล้วแก้ไขข้อความในวงเล็บเหลี่ยม)
\clearpage
\paparttitle{ส่วนที่ 3}
\padomaintitle{ผลการดำเนินงาน ด้านที่ 2 [ชื่อด้านที่ 2 ตามแบบประเมินตำแหน่ง/วิทยฐานะของผู้จัดทำ]}
\padomaintitle{ตัวชี้วัดที่ 1 [ชื่อตัวชี้วัดที่ 1]}
\vspace{6pt}

\noindent ข้าพเจ้า [คำนำหน้า ชื่อ-นามสกุล] ตำแหน่ง [ตำแหน่ง] วิทยฐานะ [วิทยฐานะ]
สังกัด [หน่วยงาน] ขอรายงานผลการดำเนินงานด้านที่ 2 ตัวชี้วัดที่ 1 ดังรายละเอียดต่อไปนี้

\noindent ข้าพเจ้าดำเนินงานภายใต้รูปแบบ [ชื่อรูปแบบ/นวัตกรรมของผู้จัดทำ] ซึ่งประกอบด้วย 8 กลยุทธ์ ดังนี้
\begin{itemize}
  \setlength\itemsep{0pt}
  \item กลยุทธ์ที่ 1  [ชื่อกลยุทธ์ที่ 1]  [คำอธิบายสั้น ๆ]
  \item กลยุทธ์ที่ 2  [ชื่อกลยุทธ์ที่ 2]  [คำอธิบายสั้น ๆ]
  \item กลยุทธ์ที่ 3  [ชื่อกลยุทธ์ที่ 3]  [คำอธิบายสั้น ๆ]
  \item กลยุทธ์ที่ 4  [ชื่อกลยุทธ์ที่ 4]  [คำอธิบายสั้น ๆ]
  \item กลยุทธ์ที่ 5  [ชื่อกลยุทธ์ที่ 5]  [คำอธิบายสั้น ๆ]
  \item กลยุทธ์ที่ 6  [ชื่อกลยุทธ์ที่ 6]  [คำอธิบายสั้น ๆ]
  \item กลยุทธ์ที่ 7  [ชื่อกลยุทธ์ที่ 7]  [คำอธิบายสั้น ๆ]
  \item กลยุทธ์ที่ 8  [ชื่อกลยุทธ์ที่ 8]  [คำอธิบายสั้น ๆ]
\end{itemize}

% ---------- เกณฑ์การให้คะแนน ----------
\paheading{เกณฑ์การให้คะแนน (Scoring Rubric) ส่วนที่ 3 ตัวชี้วัดที่ 1}
\begin{longtable}{|c|p{9.5cm}|}
\hline
\tblhead{คะแนน} & \tblhead{เกณฑ์การพิจารณาผลการปฏิบัติงาน / ผลลัพธ์} \\
\hline
\endfirsthead
\hline
\tblhead{คะแนน} & \tblhead{เกณฑ์การพิจารณาผลการปฏิบัติงาน / ผลลัพธ์} \\
\hline
\endhead
\hline\endfoot
\hline\endlastfoot
1 คะแนน & [เกณฑ์การพิจารณาระดับคะแนน 1 --- คัดลอกข้อความจากแบบประเมิน PA 4 ของ ก.ค.ศ. ตามตำแหน่ง/วิทยฐานะจริง] \\ \hline
2 คะแนน & [เกณฑ์การพิจารณาระดับคะแนน 2] \\ \hline
3 คะแนน & [เกณฑ์การพิจารณาระดับคะแนน 3] \\ \hline
4 คะแนน & [เกณฑ์การพิจารณาระดับคะแนน 4] \\ \hline
5 คะแนน & [เกณฑ์การพิจารณาระดับคะแนน 5] \\ \hline
\end{longtable}

% ---------- ตารางสรุปความสอดคล้อง ----------
\paheading{ตารางสรุป การวิเคราะห์ความสอดคล้องระหว่างเกณฑ์ PA 4 กับกลยุทธ์ของผู้จัดทำ}
\begin{longtable}{|p{3.6cm}|p{3.6cm}|p{3cm}|p{3cm}|}
\hline
\tblhead{เกณฑ์การพิจารณาผลการปฏิบัติตาม} & \tblhead{การดำเนินงานตามกลยุทธ์} & \tblhead{ผลลัพธ์} & \tblhead{หลักฐานอ้างอิง} \\
\hline
\endfirsthead
\hline
\tblhead{เกณฑ์การพิจารณาผลการปฏิบัติตาม} & \tblhead{การดำเนินงานตามกลยุทธ์} & \tblhead{ผลลัพธ์} & \tblhead{หลักฐานอ้างอิง} \\
\hline
\endhead
\hline\endfoot
\hline\endlastfoot
1) [ข้อความเกณฑ์การพิจารณาข้อที่ 1] & [การดำเนินงานที่สอดคล้อง] & [ผลลัพธ์ที่เกิดขึ้น] & [หลักฐานอ้างอิง] \\ \hline
2) [ข้อความเกณฑ์การพิจารณาข้อที่ 2] & [การดำเนินงานที่สอดคล้อง] & [ผลลัพธ์ที่เกิดขึ้น] & [หลักฐานอ้างอิง] \\ \hline
3) [ข้อความเกณฑ์การพิจารณาข้อที่ 3] & [การดำเนินงานที่สอดคล้อง] & [ผลลัพธ์ที่เกิดขึ้น] & [หลักฐานอ้างอิง] \\ \hline
4) [ข้อความเกณฑ์การพิจารณาข้อที่ 4] & [การดำเนินงานที่สอดคล้อง] & [ผลลัพธ์ที่เกิดขึ้น] & [หลักฐานอ้างอิง] \\ \hline
5) [ข้อความเกณฑ์การพิจารณาข้อที่ 5] & [การดำเนินงานที่สอดคล้อง] & [ผลลัพธ์ที่เกิดขึ้น] & [หลักฐานอ้างอิง] \\ \hline
\end{longtable}

% ---------- การดำเนินงานตามกลยุทธ์ (PDCA) ----------
\paheading{การดำเนินงานตามกลยุทธ์ที่ 1  [ชื่อกลยุทธ์]}
\noindent ข้าพเจ้า [คำนำหน้า ชื่อ-นามสกุล] ตำแหน่ง [ตำแหน่ง] วิทยฐานะ [วิทยฐานะ]
มีกระบวนการดำเนินงานเพื่อให้เกิดผลลัพธ์ตามเกณฑ์การประเมิน ดังนี้
\begin{itemize}
  \setlength\itemsep{0pt}
  \item 1.1  [รายละเอียดกิจกรรม/ผลการดำเนินงานย่อยข้อที่ 1]
  \item 1.2  [รายละเอียดกิจกรรม/ผลการดำเนินงานย่อยข้อที่ 2]
  \item 1.3  [รายละเอียดกิจกรรม/ผลการดำเนินงานย่อยข้อที่ 3]
  \item 1.4  [รายละเอียดกิจกรรม/ผลการดำเนินงานย่อยข้อที่ 4]
  \item 1.5  [รายละเอียดกิจกรรม/ผลการดำเนินงานย่อยข้อที่ 5]
\end{itemize}

\noindent ข้าพเจ้ามีกระบวนการในการดำเนินงาน ดังนี้

\pasubheading{1.  ขั้นเตรียมการ}
\noindent [บรรยายการศึกษาสภาพปัญหา/ความต้องการจำเป็น และการออกแบบกลยุทธ์/แผนงานก่อนดำเนินการ
พร้อมอ้างอิงหลักฐาน]

\pasubheading{2.  ขั้นดำเนินการ}
\noindent [บรรยายภาพรวมการนำแผนสู่การปฏิบัติ ตามวงจรคุณภาพ PDCA ดังนี้]
\begin{itemize}
  \setlength\itemsep{0pt}
  \item 2.1  ขั้นวางแผน (Plan)  [รายละเอียดขั้นวางแผน]
  \item 2.2  ขั้นปฏิบัติตามแผน (Do)  [รายละเอียดขั้นปฏิบัติ]
  \item 2.3  ขั้นตรวจสอบ (Check)  [รายละเอียดขั้นตรวจสอบ/ติดตาม/ประเมินผล]
  \item 2.4  ขั้นพัฒนาปรับปรุง (Act)  [รายละเอียดการนำผลประเมินไปปรับปรุง]
\end{itemize}

\pasubheading{3.  ขั้นจัดทำรายงานประเมินผล}
\noindent [บรรยายการจัดทำรายงานสรุปผลการดำเนินงาน และการนำเสนอ/เผยแพร่ผลงาน]

% ---------- ตารางที่ 1 ----------
\pasubheading{ตารางที่ 1  การวิเคราะห์ความสอดคล้องกับเกณฑ์การให้คะแนน ตัวชี้วัดที่ 1}
\begin{longtable}{|c|p{5cm}|p{6.5cm}|}
\hline
\tblhead{ข้อ} & \tblhead{เกณฑ์การพิจารณาตาม PA 4} & \tblhead{ผลการปฏิบัติงานและหลักฐานเชิงประจักษ์} \\
\hline
\endfirsthead
\hline
\tblhead{ข้อ} & \tblhead{เกณฑ์การพิจารณาตาม PA 4} & \tblhead{ผลการปฏิบัติงานและหลักฐานเชิงประจักษ์} \\
\hline
\endhead
\hline\endfoot
\hline\endlastfoot
1 & [เกณฑ์การพิจารณาตาม PA 4 ข้อที่ 1] & [ผลการปฏิบัติงานและหลักฐานเชิงประจักษ์ข้อที่ 1] \\ \hline
2 & [เกณฑ์การพิจารณาตาม PA 4 ข้อที่ 2] & [ผลการปฏิบัติงานและหลักฐานเชิงประจักษ์ข้อที่ 2] \\ \hline
3 & [เกณฑ์การพิจารณาตาม PA 4 ข้อที่ 3] & [ผลการปฏิบัติงานและหลักฐานเชิงประจักษ์ข้อที่ 3] \\ \hline
4 & [เกณฑ์การพิจารณาตาม PA 4 ข้อที่ 4] & [ผลการปฏิบัติงานและหลักฐานเชิงประจักษ์ข้อที่ 4] \\ \hline
5 & [เกณฑ์การพิจารณาตาม PA 4 ข้อที่ 5] & [ผลการปฏิบัติงานและหลักฐานเชิงประจักษ์ข้อที่ 5] \\ \hline
\end{longtable}

% ---------- ผลลัพธ์ + ตารางที่ 2 ----------
\paheading{ผลลัพธ์}
\noindent [สรุปผลลัพธ์ภาพรวมที่เกิดขึ้นจากการดำเนินงานตามกลยุทธ์นี้ พร้อมตัวเลข/ร้อยละสนับสนุน]

\pasubheading{ตารางที่ 2  ผลการดำเนินงาน (ตัวอย่างรูปแบบตาราง --- แก้ไขหัวข้อ/จำนวนแถวตามจริง)}
\padisclaimer
\begin{longtable}{|c|p{3cm}|p{2.6cm}|p{1.8cm}|p{1.3cm}|p{3.2cm}|}
\hline
\tblhead{ที่} & \tblhead{กิจกรรม} & \tblhead{กลุ่มเป้าหมาย} & \tblhead{เข้าร่วม} & \tblhead{ร้อยละ} & \tblhead{ผลการดำเนินงาน} \\
\hline
\endfirsthead
\hline
\tblhead{ที่} & \tblhead{กิจกรรม} & \tblhead{กลุ่มเป้าหมาย} & \tblhead{เข้าร่วม} & \tblhead{ร้อยละ} & \tblhead{ผลการดำเนินงาน} \\
\hline
\endhead
\hline\endfoot
\hline\endlastfoot
1 & [กิจกรรมที่ 1] & [กลุ่มเป้าหมาย] & [จำนวน] & 0.00 & [ผลการดำเนินงาน] \\ \hline
2 & [กิจกรรมที่ 2] & [กลุ่มเป้าหมาย] & [จำนวน] & 0.00 & [ผลการดำเนินงาน] \\ \hline
3 & [กิจกรรมที่ 3] & [กลุ่มเป้าหมาย] & [จำนวน] & 0.00 & [ผลการดำเนินงาน] \\ \hline
4 & [กิจกรรมที่ 4] & [กลุ่มเป้าหมาย] & [จำนวน] & 0.00 & [ผลการดำเนินงาน] \\ \hline
5 & [กิจกรรมที่ 5] & [กลุ่มเป้าหมาย] & [จำนวน] & 0.00 & [ผลการดำเนินงาน] \\ \hline
\end{longtable}
\noindent จากตารางที่ 2 พบว่า [สรุปผลภาพรวมจากตาราง]

% ---------- รายการหลักฐานเชิงประจักษ์รายข้อ ----------
\pasubheading{1.1.1  [ชื่อกิจกรรม/ผลลัพธ์ย่อยข้อที่ 1]}
\noindent [บรรยายเชื่อมโยงว่ากิจกรรม/ผลลัพธ์ข้อที่ 1 สอดคล้องกับตัวชี้วัดที่ 1 อย่างไร]\\
\textbf{ผลการดำเนินงาน :} [ระบุผลเชิงประจักษ์/ตัวเลขยืนยันของข้อที่ 1]
\paphotobox

\pasubheading{1.1.2  [ชื่อกิจกรรม/ผลลัพธ์ย่อยข้อที่ 2]}
\noindent [บรรยายเชื่อมโยงว่ากิจกรรม/ผลลัพธ์ข้อที่ 2 สอดคล้องกับตัวชี้วัดที่ 1 อย่างไร]\\
\textbf{ผลการดำเนินงาน :} [ระบุผลเชิงประจักษ์/ตัวเลขยืนยันของข้อที่ 2]
\paphotobox

\pasubheading{1.1.3  [ชื่อกิจกรรม/ผลลัพธ์ย่อยข้อที่ 3]}
\noindent [บรรยายเชื่อมโยงว่ากิจกรรม/ผลลัพธ์ข้อที่ 3 สอดคล้องกับตัวชี้วัดที่ 1 อย่างไร]\\
\textbf{ผลการดำเนินงาน :} [ระบุผลเชิงประจักษ์/ตัวเลขยืนยันของข้อที่ 3]
\paphotobox

\pasubheading{1.1.4  [ชื่อกิจกรรม/ผลลัพธ์ย่อยข้อที่ 4]}
\noindent [บรรยายเชื่อมโยงว่ากิจกรรม/ผลลัพธ์ข้อที่ 4 สอดคล้องกับตัวชี้วัดที่ 1 อย่างไร]\\
\textbf{ผลการดำเนินงาน :} [ระบุผลเชิงประจักษ์/ตัวเลขยืนยันของข้อที่ 4]
\paphotobox

\pasubheading{1.1.5  [ชื่อกิจกรรม/ผลลัพธ์ย่อยข้อที่ 5]}
\noindent [บรรยายเชื่อมโยงว่ากิจกรรม/ผลลัพธ์ข้อที่ 5 สอดคล้องกับตัวชี้วัดที่ 1 อย่างไร]\\
\textbf{ผลการดำเนินงาน :} [ระบุผลเชิงประจักษ์/ตัวเลขยืนยันของข้อที่ 5]
\paphotobox

% ---------- ตารางที่ 3 ----------
\pasubheading{ตารางที่ 3  สรุปผลการดำเนินงานเทียบค่าเป้าหมาย}
\padisclaimer
\begin{longtable}{|p{4.3cm}|p{2cm}|p{2cm}|p{2cm}|p{2cm}|}
\hline
\tblhead{ประเด็นการพิจารณา} & \tblhead{ข้อมูลฐาน (ร้อยละ)} & \tblhead{ค่าเป้าหมาย (ร้อยละ)} & \tblhead{ผลการประเมิน (ร้อยละ)} & \tblhead{สรุปผล} \\
\hline
\endfirsthead
\hline
\tblhead{ประเด็นการพิจารณา} & \tblhead{ข้อมูลฐาน (ร้อยละ)} & \tblhead{ค่าเป้าหมาย (ร้อยละ)} & \tblhead{ผลการประเมิน (ร้อยละ)} & \tblhead{สรุปผล} \\
\hline
\endhead
\hline\endfoot
\hline\endlastfoot
1) [ประเด็นการพิจารณาข้อที่ 1] & 0.00 & [เป้าหมาย] & 0.00 & [บรรลุ/ไม่บรรลุ] \\ \hline
2) [ประเด็นการพิจารณาข้อที่ 2] & 0.00 & [เป้าหมาย] & 0.00 & [บรรลุ/ไม่บรรลุ] \\ \hline
3) [ประเด็นการพิจารณาข้อที่ 3] & 0.00 & [เป้าหมาย] & 0.00 & [บรรลุ/ไม่บรรลุ] \\ \hline
4) [ประเด็นการพิจารณาข้อที่ 4] & 0.00 & [เป้าหมาย] & 0.00 & [บรรลุ/ไม่บรรลุ] \\ \hline
5) [ประเด็นการพิจารณาข้อที่ 5] & 0.00 & [เป้าหมาย] & 0.00 & [บรรลุ/ไม่บรรลุ] \\ \hline
\end{longtable}
\noindent จากตารางที่ 3 พบว่า [สรุปผลภาพรวมเทียบค่าเป้าหมาย]

\noindent ข้อสังเกตที่แสดงถึงระดับการปฏิบัติที่คาดหวังของวิทยฐานะที่ขอรับการประเมิน คือ
[ระบุข้อสังเกต/จุดเด่นที่ยกระดับจากวิทยฐานะเดิม]

% ---------- ตารางแหล่งข้อมูลสำหรับตรวจสอบ (ลบก่อนยื่นจริง) ----------
\pasubheading{ตารางแหล่งข้อมูลสำหรับตรวจสอบ / แทนที่ตัวเลขสมมุติ (สำหรับผู้จัดทำ --- ลบตารางนี้ออกก่อนยื่นเอกสารจริง)}
\noindent ตารางนี้เป็นแนวทางบอกว่าแต่ละช่องตัวเลขในตารางที่ 2 และตารางที่ 3 ควรดึงข้อมูลจากเอกสารหรือระบบใด
เพื่อให้ตัวเลขที่กรอกตรวจสอบย้อนกลับได้
\begin{longtable}{|c|p{5.5cm}|p{6cm}|}
\hline
\tblhead{รหัส} & \tblhead{แหล่งข้อมูล} & \tblhead{วิธีนับ / หมายเหตุ} \\
\hline
\endfirsthead
\hline
\tblhead{รหัส} & \tblhead{แหล่งข้อมูล} & \tblhead{วิธีนับ / หมายเหตุ} \\
\hline
\endhead
\hline\endfoot
\hline\endlastfoot
ก & [ชื่อเอกสาร/ระบบที่เป็นแหล่งข้อมูล] & [วิธีนับ/หมายเหตุ] \\ \hline
ข & [ชื่อเอกสาร/ระบบที่เป็นแหล่งข้อมูล] & [วิธีนับ/หมายเหตุ] \\ \hline
ค & [ชื่อเอกสาร/ระบบที่เป็นแหล่งข้อมูล] & [วิธีนับ/หมายเหตุ] \\ \hline
ง & [ชื่อเอกสาร/ระบบที่เป็นแหล่งข้อมูล] & [วิธีนับ/หมายเหตุ] \\ \hline
จ & [ชื่อเอกสาร/ระบบที่เป็นแหล่งข้อมูล] & [วิธีนับ/หมายเหตุ] \\ \hline
\end{longtable}

% ---------- ลงชื่อผู้รายงาน ----------
\pasignatureblock{( [คำนำหน้า ชื่อ-นามสกุล] )}{ตำแหน่ง [ตำแหน่ง]  วิทยฐานะ [วิทยฐานะ]}{[หน่วยงาน/สำนักงานเขตพื้นที่การศึกษา]}
